\documentclass[11pt]{article}

\usepackage[margin=1in]{geometry}
\usepackage{amsmath, amssymb}
\usepackage{booktabs}
\usepackage{graphicx}
\usepackage{lmodern}      % scalable fonts; microtype's expansion needs them
\usepackage[T1]{fontenc}
\usepackage{microtype}
\usepackage{xcolor}
\usepackage[colorlinks=true, allcolors=blue!60!black]{hyperref}
\usepackage{caption}
\captionsetup{font=small, labelfont=bf}

\title{\textbf{What Transfers and What Does Not:\\
Reproducing Compressed Sparse Attention at Single-GPU Scale}}
\author{Brett Cleary}
\date{September 2026}

\begin{document}
\maketitle

\begin{abstract}
We attempted to reproduce the central claim of Compressed Sparse Attention
(CSA), the attention mechanism of DeepSeek-V4, at a scale reachable by one
consumer GPU: roughly $10^7$ parameters, byte-level modelling, and context
lengths from 1K to 64K. We did not succeed, and the reasons are more
informative than the intended result would have been.

Three findings concern measurement rather than architecture. First, the
published equations under-determine the architecture: implementing
Eqs.~(9)--(19) faithfully omits a sliding-window branch that appears only in
the accompanying figure, leaving every query blind to the $t \bmod m$ bytes
immediately preceding it, at a cost of roughly $0.52$ bits per character.
Second, our dense baseline lacked the query--key normalisation that CSA carries
by specification; with a gradient-norm clip of $1.0$, this left the baseline
clipped on $94\%$ of steps at 16K context and training at approximately
one-fifteenth of its nominal learning rate, an artefact larger than the
architecture effect under study. Third, mean bits-per-character---the standard
metric---dilutes the long-context effect by roughly $80\times$: the median byte
gains $0.0002$ bpc from an additional 8K of context while the top $1\%$ of
bytes gain $1.83$ bpc and account for $81\%$ of the total.

The fourth finding concerns the efficiency claim itself, and splits it in two.
CSA's compute advantage did not survive reimplementation: against a dense
baseline dispatching to a FlashAttention kernel, our CSA used $4.3\times$ the
memory and ran $3.8\times$ slower at 8K, and exhausted memory at 16K on
hardware where dense trained comfortably at 64K. Its KV-cache advantage, by
contrast, reproduced exactly---$32\times$ at our configuration, invariant in
context length---because it is a property of the architecture's shape rather
than of its kernels. We argue this distinction is the transferable lesson:
architectural claims survive small-scale reproduction, engineering-dependent
claims do not, and published descriptions rarely make clear which is which.
\end{abstract}

\section{Introduction}

Architecture research at frontier scale is not reproducible by independent
researchers. A 1.6T-parameter model trained on 33T tokens is beyond
verification outside a handful of laboratories. The implicit hope is that the
\emph{mechanisms} such work introduces are separable from the scale at which
they were demonstrated: that one can implement an attention variant on a single
GPU, observe a directionally similar benefit, and thereby verify the idea even
if not the artefact.

This paper tests that hope on a specific case. Compressed Sparse Attention
(CSA), introduced in the DeepSeek-V4 technical report~\cite{deepseekv4},
compresses the key--value cache by a factor $m$ along the sequence dimension
and then applies sparse top-$k$ selection over the compressed entries. It is
presented as the enabling mechanism for million-token context. We implemented
it from the paper's equations, validated the implementation with unit tests,
and set out to answer a single question:

\begin{quote}
\emph{Does compressed sparse attention overtake dense attention as context
length grows, at a scale a single GPU can reach?}
\end{quote}

The answer is that the question is malformed, and discovering how it is
malformed took four distinct experimental failures. We report them because each
failure is a concrete, quantified obstacle that any small-scale reproduction of
a long-context attention mechanism will encounter, and because three of the
four were invisible without instrumentation we did not initially think to
build.

\paragraph{Contributions.}
\begin{enumerate}
  \item We show the published CSA equations under-determine the architecture,
        and quantify the cost of the resulting omission
        (Section~\ref{sec:recency}).
  \item We identify and measure a baseline-construction artefact---missing
        query--key normalisation---that exceeded the architecture effect being
        studied, and that grows with context length
        (Section~\ref{sec:qknorm}).
  \item We show that mean bits-per-character is the wrong instrument for
        long-context evaluation, and give a corruption-based measurement that
        recovers the effect it hides (Section~\ref{sec:dilution}).
  \item We show that CSA's compute advantage inverts without custom kernels
        while its KV-cache advantage reproduces exactly, and argue this is the
        general pattern (Section~\ref{sec:efficiency}).
\end{enumerate}

We explicitly do \emph{not} report a clean quality comparison between CSA and
dense attention. Section~\ref{sec:sweep} explains why the comparisons we ran
are uninterpretable, and we consider publishing them as though they were not to
be the more serious error.

\section{Background}

\subsection{Compressed Sparse Attention}

CSA (\cite{deepseekv4}, \S2.3.1) proceeds in three stages. Given hidden states
$H \in \mathbb{R}^{n \times d}$, it first forms two series of KV entries
$C^a, C^b \in \mathbb{R}^{n \times c}$ with compression weights
$Z^a, Z^b$, and compresses each group of $m$ tokens into a single entry
$C^{\mathrm{Comp}}_i \in \mathbb{R}^{c}$ by a softmax-weighted sum over a
$2m$-token window (Eqs.~9--12). Adjacent compressed entries overlap in the
token ranges they draw from, so the sequence is compressed to $n/m$ entries.

A \emph{lightning indexer} then scores each compressed block against a query
token through a low-rank query path shared with the core attention (Eqs.~13--17),
and a top-$k$ selector retains only the highest-scoring blocks. Finally, core
attention is performed in multi-query fashion~\cite{shazeer2019mqa}, with every
query head attending to the same retained compressed entries, which serve as
both keys and values (Eqs.~18--19).

Two details matter for what follows. The paper applies
RMSNorm~\cite{zhang2019rmsnorm} to queries and to compressed KV entries,
stating that this ``avoids exploding attention logits and may improve training
stability'' (\S2.3.3). And it introduces an additional sliding-window attention
branch, whose motivation we return to in Section~\ref{sec:recency}.

\subsection{What CSA claims}

It is worth stating precisely, because it constrains what a reproduction can
show. CSA is introduced ``to enhance long-context \emph{efficiency}''; the
claim is cost, not quality. The training recipe makes this concrete: DeepSeek-V4
warms up with \emph{dense} attention for the first 1T tokens and introduces
sparse attention only at a sequence length of 64K, having trained at 4K and 16K
densely (\S4.2.2). CSA is never trained from scratch and never used below 64K.

This has a consequence that took us some time to internalise. Compression
($m$ tokens to one entry) and selection (top-$k$ of those) are both lossy.
At matched parameters, tokens and context, dense attention is an upper bound on
what CSA can extract. Expecting CSA to \emph{surpass} dense in bits-per-character
is expecting a lossy approximation to beat what it approximates. The achievable
target is parity at lower cost, and our framing question, quoted in
Section~1, was malformed from the start in asking for more.

\section{Experimental setup}

All experiments use a decoder-only transformer with pre-norm
RMSNorm~\cite{zhang2019rmsnorm}, SwiGLU feed-forward blocks, learned absolute
position embeddings, and a tied language-model head. Unless stated otherwise:
$d_{\text{model}}=384$, 6 layers, 6 heads (head dimension 64), giving 10.6M
non-embedding parameters for the dense arm and 10.0M for CSA. CSA uses $m=4$,
$c=64$, and $\text{top-}k = n_{\text{blk}}/4$ (constant 25\% sparsity).
Optimisation is AdamW~\cite{loshchilov2019adamw} at $3\times10^{-4}$ decayed
cosine to $3\times10^{-5}$, weight decay $0.1$, gradient-norm clip $1.0$, bf16
autocast, effective batch 32 sequences. Hardware is a single RTX 5070 Ti
(16GB) for the measurements in Sections~\ref{sec:dilution} and
\ref{sec:efficiency}; the sweeps in Section~\ref{sec:sweep} used rented H100
and A100 instances.

Data is byte-level. We use enwik8~\cite{hutter2006} (90M/5M/5M
train/validation/test) and PG-19~\cite{rae2019pg19}, the standard long-range
language-modelling benchmark, subsampled to the same split sizes. All results
are bits per character (equivalently bits per byte),
$\mathrm{bpc} = \mathcal{L}/\ln 2$, on the held-out test split, evaluated from
the best-validation checkpoint. For PG-19, evaluation and training windows are
document-aware: no window crosses a book boundary. This matters---on enwik8 at
16K context, 41\% of corpus bytes lie in an article shorter than the window,
against 0.1\% for PG-19.

\section{The intended experiment, and three reasons it failed}
\label{sec:sweep}

Our intended headline experiment was a sweep over context length
$\{1\text{K}, 2\text{K}, 4\text{K}, 8\text{K}\}$ crossed with
$\{\text{dense}, \text{CSA}\}$ at fixed parameters and token budget, looking
for the context length at which the arms cross. We ran it three times.
Table~\ref{tab:sweeps} shows the result.

\begin{table}[t]
\centering
\small
\begin{tabular}{lcccccc}
\toprule
& \multicolumn{2}{c}{sweep v2} & \multicolumn{2}{c}{sweep v3} & \multicolumn{2}{c}{sweep v3.1} \\
\cmidrule(lr){2-3}\cmidrule(lr){4-5}\cmidrule(lr){6-7}
context & dense & CSA & dense & CSA & dense & CSA \\
\midrule
1K & 1.375 & 1.968 & 1.338 & 1.903 & --- & --- \\
2K & 1.420 & 2.054 & 1.352 & 1.927 & 1.419 & 1.836 \\
4K & 1.565 & 2.443 & 1.506 & 2.048 & 1.643 & 1.934 \\
8K & 2.250 & 3.180 & 1.625 & 2.691 & 1.769 & 1.900 \\
\midrule
gap at 1K & \multicolumn{2}{c}{0.593} & \multicolumn{2}{c}{0.565} & \multicolumn{2}{c}{---} \\
gap at 8K & \multicolumn{2}{c}{0.930} & \multicolumn{2}{c}{1.066} & \multicolumn{2}{c}{0.131} \\
\bottomrule
\end{tabular}
\caption{Test bpc for the same grid, run three times. The architecture and
data are identical; only the iteration caps and stopping behaviour changed.
The gap widens with context (v2), is flat then jumps (v3), or collapses toward
crossover (v3.1). None of these are trustworthy; see
Sections~\ref{sec:lrhorizon}--\ref{sec:recency}.}
\label{tab:sweeps}
\end{table}

Three sweeps produced three incompatible stories from identical architecture and
data. We diagnosed three separate causes, and none of them is about attention.

\subsection{The learning-rate schedule was tied to the stopping criterion}
\label{sec:lrhorizon}

Our cosine schedule annealed over \texttt{max\_iters}, while early stopping
halted runs before reaching it. A run that stops early therefore never completes
its anneal, and two cells that stop at different fractions of their cap are not
comparable. The cleanest demonstration is a pair of runs differing only in cap:
identical seed, context, batch, and identical best step (2500), yet
$1.625$ bpc against $1.769$ bpc---a difference of $0.144$ bpc attributable
entirely to the cosine horizon.

The bias was systematic rather than random. In sweep v3.1, every dense cell
stopped at 20--31\% of its schedule while every CSA cell ran to 51--93\% of
its own. The apparent gap collapse at 8K is what that asymmetry looks like.
We now set the anneal horizon explicitly, independent of the stopping cap.

\subsection{Only one arm had query--key normalisation}
\label{sec:qknorm}

The more serious error was in the baseline. Our CSA implementation applied
RMSNorm to queries and compressed keys, because the paper specifies it. Our
dense baseline applied none, going directly from the QKV projection into
scaled dot-product attention. We had been careful to give both arms identical
positional encoding for fairness; we did not think to check normalisation.

The consequence is Table~\ref{tab:qknorm}. Gradient norms are recorded
pre-clip, against a clip threshold of $1.0$.

\begin{table}[t]
\centering
\small
\begin{tabular}{llccc}
\toprule
arm & context & median grad.\ norm & steps clipped & effective LR \\
\midrule
dense & 1K  & 0.52 & 10.8\% & $\times 1.00$ \\
dense & 2K  & 0.77 & 41.5\% & $\times 1.00$ \\
dense & 8K  & 10.69 & \textbf{91.0\%} & $\mathbf{\times 0.094}$ \\
dense & 16K & 15.42 & \textbf{94.0\%} & $\mathbf{\times 0.065}$ \\
\midrule
CSA & 1K & 0.58 & 0.3\%  & $\times 1.00$ \\
CSA & 2K & 0.52 & 0.8\%  & $\times 1.00$ \\
CSA & 8K & 0.62 & 17.6\% & $\times 1.00$ \\
\bottomrule
\end{tabular}
\caption{Gradient norms by arm and context. The dense baseline's gradients grow
with context; CSA's do not, because CSA normalises queries and keys by
specification. At 8K and 16K the baseline takes steps roughly one tenth of
their nominal size while CSA runs essentially unclipped.}
\label{tab:qknorm}
\end{table}

This artefact is both arm-asymmetric and context-dependent, so it bites hardest
exactly where the claim under test lives. It plausibly accounts for the dense
arm early-stopping so early at long context, for dense bpc \emph{worsening}
with context in all three sweeps, and for much of v3.1's apparent gap collapse.

It became undeniable when a 16K run on PG-19 diverged outright: gradient norms
grew from $0.3$ to $4.3\times10^{5}$, best validation stranded at step 2100 of
10000, with everything after degrading by roughly 50\%. Training and validation
bpc moved together throughout, ruling out overfitting. Adding QK-normalisation
to the baseline---standard practice for dense attention~\cite{henry2020qknorm,
dehghani2023vit22b}, and what CSA already had---held gradient norms flat at
$0.33$--$0.50$ for all 10000 steps and improved best validation from
$2.1876$ to $1.4918$ bpc.

That single normalisation layer was worth $0.70$ bpc in the baseline. The
architecture effect the study set out to measure spans $0.13$--$1.07$ bpc.
The baseline was the experiment.

\subsection{The equations omit a component the figure includes}
\label{sec:recency}

CSA's causal mask permits a query at position $t$ to attend only to compressed
blocks covering positions $[0, m\lfloor t/m \rfloor)$. The $t \bmod m$ tokens
immediately preceding $x_t$ are therefore unreachable through attention at
\emph{every} layer. We verified this with a gradient probe: for $m=4$,
$\partial\,\mathrm{logits}[23] / \partial\,\mathrm{embedding}[s]$ is exactly
zero for $s \in \{20,21,22\}$ at any depth, while the dense arm has no such
gaps.

This is not a bug in our implementation of the equations. It is a property of
the compressed branch, and the paper says so directly (\S2.3.3):

\begin{quote}\small
``In order to strictly preserve causality in CSA and HCA, each query attends to
only preceding compressed KV blocks. Consequently, a query cannot access
information from other tokens within its own compressed block. \dots For these
reasons, we introduce a supplementary attention branch \dots in a sliding
window manner.''
\end{quote}

The remedy is a sliding-window branch of $n_{\text{win}}$ uncompressed entries,
concatenated with the selected compressed entries before core attention. It
appears in the CSA architecture figure, on the same page as the equations, as a
concatenation node feeding the shared-KV attention. It does not appear in
Eqs.~(9)--(19): Eq.~(19) takes keys and values to be the selected compressed
entries alone.

An implementation built faithfully from the equations is therefore missing a
component that the figure includes and that the text motivates explicitly. An
order-$K$ byte $n$-gram proxy prices the resulting blind spot at
$+0.514$ bpc averaged over $t \bmod m$ (stable at $+0.521$ for $K=8$), which is
comparable to the entire dense--CSA gap our sweeps reported at short context.
We have since implemented the branch ($n_{\text{win}}$ configurable; DeepSeek-V4
uses 128 at $m=4$).

We note that at our scale the paper's constant is awkward: $n_{\text{win}}=128$
is $0.013\%$ of a 1M context but $12.5\%$ of a 1K context. A fixed architectural
constant tuned for one regime does not transfer to another three orders of
magnitude away, and at 1K it would supply a dense local window over an eighth of
the sequence.

\subsection{Consequence}

The three artefacts do not share a direction. The learning-rate horizon and the
QK-normalisation asymmetry both flatter CSA, and both worsen with context. The
recency hole flatters dense. Each is comparable to or larger than the effect
being measured. We therefore report no corrected version of
Table~\ref{tab:sweeps}: the correct conclusion is that these measurements are
uninterpretable, not that they can be adjusted.

\section{Mean bits-per-character is the wrong instrument}
\label{sec:dilution}

Having fixed the baseline, we asked a prior question: does context length
change the outcome enough to be worth sweeping over? We trained a dense 16K
model on PG-19 to completion (gradient norms flat, anneal completed, best
validation at step 9700 of 10000) and measured test bpc bucketed by position
within the evaluation window, since a token at position $t$ has $t$ bytes of
available context.

The curve is nearly flat: $1.654$ bpc over positions 0--1024, $1.546$ over
1024--2048, then drifting slowly to $1.513$ at 16K. A paired per-window test
between an early slice (2--4K of context) and a late slice (14--16K) gives
$+0.0240 \pm 0.0086$ bpc ($t=2.8$, 326 windows). Statistically detectable;
practically negligible against architecture effects of $0.13$--$1.07$ bpc.

We initially concluded that the context axis was inert and that the study was
not viable on this data. That conclusion was wrong, and the error is
instructive: mean bpc averages over every position, while long-range dependence
affects only a small subset of them. If $1\%$ of bytes gain 2 bits and $99\%$
gain nothing, the mean gain is $0.02$ bpc---indistinguishable from no effect.

To measure the distribution instead, we score each token twice on sequences
identical in length, absolute position, and local context, differing only in
whether the \emph{distant} context carries real information: bytes before a
cutoff are replaced with prose from a different book. Table~\ref{tab:dilution}
gives the result.

\begin{table}[t]
\centering
\small
\begin{tabular}{lr}
\toprule
statistic & value \\
\midrule
mean gain from true distant context & $+0.0226$ bpc \\
\textbf{median gain} & $\mathbf{+0.0002}$ \textbf{bpc} \\
tokens helped / hurt & 54.2\% / 45.8\% \\
tokens gaining $>1$ bit & 0.96\%, carrying 79\% of all bits gained \\
top 1\% of tokens & $+1.83$ bpc each, 81\% of all bits gained \\
\bottomrule
\end{tabular}
\caption{Value of genuine distant context, PG-19, 327{,}680 scored positions.
Context before byte 8192 replaced with text from another book; only positions
12288 and beyond scored, so every scored token retains at least 4K of genuine
local context. The mean agrees with the independent position-bucketed
measurement ($+0.024$); the median does not.}
\label{tab:dilution}
\end{table}

The typical byte gains nothing from 8K of additional context---it is close to a
coin flip whether it gains or loses. The effect is concentrated in roughly $1\%$
of tokens gaining $1.83$ bits each. Averaging dilutes a large effect by
approximately $80\times$.

This reframes the evaluation problem. On the subpopulation that depends on
distant retrieval, the effect is $1.83$ bpc: larger than every artefact
catalogued in Section~\ref{sec:sweep}, rather than far smaller. A long-context
architecture comparison at this scale is feasible, but only conditioned on the
tokens where long context matters.

\paragraph{Two designs that failed.} Both produced apparent gains above
4 bpc---larger than the model's entire bits-per-character, which is the
signature to watch for. Masking attention to a short window at evaluation takes
a model trained with full attention off its training distribution. Re-feeding
only the recent bytes as their own sequence is worse: the sequence then
contains no position-0 token, and transformers depend on an early-token
attention sink~\cite{xiao2024streamingllm}, so the model collapses to near-random
($5.70$ bpc against $6.98$ for a uniform guess over this vocabulary).
Substituting real text from another document preserves length, positions, and
the sink, and changes only information content.

\section{The efficiency claim splits in two}
\label{sec:efficiency}

CSA's claim is efficiency, so we measured it directly rather than inferring it
from FLOP counts. Table~\ref{tab:feasibility} reports peak memory and
throughput for a forward, backward, and optimiser step at batch size 1.

\begin{table}[t]
\centering
\small
\begin{tabular}{lcccc}
\toprule
& \multicolumn{2}{c}{dense (FlashAttention SDPA)} & \multicolumn{2}{c}{CSA ($m=4$, 25\% top-$k$)} \\
\cmidrule(lr){2-3}\cmidrule(lr){4-5}
context & peak memory & throughput & peak memory & throughput \\
\midrule
8K  & 1.53 GiB  & 251K tok/s & 6.65 GiB & 66K tok/s \\
16K & 2.91 GiB  & 176K tok/s & \multicolumn{2}{c}{out of memory} \\
32K & 5.62 GiB  & 110K tok/s & \multicolumn{2}{c}{out of memory} \\
64K & \textbf{11.06 GiB} & \textbf{67K tok/s} & \multicolumn{2}{c}{out of memory} \\
\bottomrule
\end{tabular}
\caption{Training feasibility on one RTX 5070 Ti (13.6 GiB free). Dense trains
at 64K context; our CSA cannot reach 16K, and at 8K uses $4.3\times$ the memory
while running $3.8\times$ slower.}
\label{tab:feasibility}
\end{table}

The result inverts the premise. On this hardware CSA is strictly worse than
dense on memory, on speed, and on maximum reachable context. The reason is that
the comparison is not FLOPs against FLOPs. Dense attention dispatches to a
FlashAttention kernel~\cite{dao2022flashattention}, which is $O(N)$ in memory
and never materialises the $N \times N$ score matrix. Our CSA materialises an
$(b, n_h, n, n_{\text{blk}})$ score tensor, and autocast computes softmax in
fp32, so a single 16K layer holds 1.5GB of attention probabilities. CSA's
$4\times$ saving in attention FLOPs is swamped by roughly a $15\times$ gap in
kernel quality.

This is not a claim that CSA is inefficient. It is a claim that CSA's
efficiency is a property of the kernels DeepSeek built for it---the report
describes a domain-specific language and mixed FP8/FP4 storage for exactly this
purpose---and that those kernels, not the equations, carry the advantage. A
reimplementation of the mathematics reproduces none of it.

\subsection{What does reproduce}

The KV-cache claim behaves differently. Dense attention must store $K$ and $V$
for every past token and every head: $2 n_h d_h$ values per token per layer.
CSA stores one compressed entry per $m$ tokens, shared across all heads by
multi-query attention, plus the indexer keys. For our configuration this is
$9216$ bytes per token against $288$, a $32\times$ reduction, invariant in
context length (Table~\ref{tab:kvcache}).

\begin{table}[h]
\centering
\small
\begin{tabular}{lrrr}
\toprule
context & dense KV cache & CSA KV cache & reduction \\
\midrule
16K  & 144.0 MB & 4.6 MB & $31\times$ \\
64K  & 576.0 MB & 18.1 MB & $32\times$ \\
256K & 2304.0 MB & 72.1 MB & $32\times$ \\
1M   & 9216.0 MB & 288.1 MB & $32\times$ \\
\bottomrule
\end{tabular}
\caption{KV cache at inference. FlashAttention does not help dense here: it
avoids materialising attention scores during compute, but autoregressive
decoding must still store every past key and value.}
\label{tab:kvcache}
\end{table}

This number requires no kernel engineering, no training, and no scale. It
follows from head count, $m$, $c$, and multi-query sharing, and it holds
identically for a 10M-parameter model on a consumer GPU and for a 1.6T-parameter
model on a cluster. It is consistent in kind with the report's claim of
approximately 2\% of a GQA-8 baseline, which additionally interleaves a more
aggressively compressed attention variant.

\section{Discussion}

\paragraph{Architectural claims transfer; engineering claims do not.}
The two halves of Section~\ref{sec:efficiency} are the paper's central
observation. CSA's KV-cache reduction is arithmetic on the shape of the
architecture, and it reproduced exactly at $10^{-5}$ of the original scale.
CSA's throughput advantage is a property of fused kernels, and it did not merely
fail to reproduce---it inverted. Both are presented in the source report as
aspects of a single efficiency story. Nothing in the presentation distinguishes
the claim a reader can verify with a calculator from the claim that requires a
compiler team.

We suspect this generalises. Sparsity, compression, and cache-sharing
mechanisms are typically justified by asymptotic FLOP or memory arguments, while
their realised benefit depends on whether a competitive kernel exists. A dense
baseline in 2026 has FlashAttention; a novel sparse mechanism usually does not,
outside the laboratory that introduced it. Small-scale reproduction will
systematically understate such mechanisms unless the reproduction also
reimplements the kernels, which is a different project with a different skill
set.

\paragraph{The baseline is the experiment.}
Our largest measured artefact was not in the mechanism under study but in what
we compared it against. A missing normalisation layer, standard in contemporary
dense attention and specified for CSA, was worth $0.70$ bpc and a
$15\times$ difference in effective learning rate at 16K context---exceeding the
architecture effect. It was invisible in loss curves and surfaced only from
logged gradient norms. We would now regard per-arm gradient-norm statistics and
clip-rate statistics as mandatory reporting for any architecture comparison, on
the grounds that a gradient clip converts an instability difference into a
silent learning-rate difference.

\paragraph{Specification by figure.}
The sliding-window branch is specified in a figure and a prose paragraph but
absent from the numbered equations. An implementer working from the equations,
which is the natural reading order for a mechanism given in nineteen numbered
steps, produces a model with a structural blind spot at every layer. We do not
think this is unusual, and we suggest that reports of new mechanisms state
explicitly which components are load-bearing for the reported results.

\paragraph{Metric selection.}
Mean bits-per-character is standard for character-level language modelling and
is the wrong instrument for long-context questions. It dilutes the phenomenon by
roughly $80\times$ in our measurement. Any evaluation of a long-context
mechanism that reports only aggregate perplexity or bpc is, on this evidence,
underpowered by one to two orders of magnitude.

\section{Limitations}

We do not report a valid quality comparison between CSA and dense attention.
The sweeps in Table~\ref{tab:sweeps} predate all three corrections; a corrected
sweep, conditioned on the subpopulation identified in
Section~\ref{sec:dilution}, remains to be run. All results use a single seed.
Our models are approximately 10M non-embedding parameters at byte-level
granularity with learned absolute position embeddings rather than rotary
embeddings, so conclusions about what a model can extract from distant context
are specific to that regime; a larger model, or one with better positional
generalisation, may exploit long range differently.

The efficiency inversion in Table~\ref{tab:feasibility} is a statement about a
straightforward PyTorch implementation, not about CSA. We did not attempt a
fused kernel, and a competent one would change the comparison substantially. We
regard establishing \emph{how much} of the gap is recoverable without
laboratory-scale engineering as the most valuable follow-up this work suggests.

\section{Conclusion}

We set out to find the smallest scale at which compressed sparse attention
demonstrates a measurable advantage over dense attention, and instead found
that at single-GPU scale the measurement is dominated by artefacts of the
comparison rather than by the mechanism. A missing baseline normalisation
exceeded the effect under study; a component present only in a figure cost
roughly $0.52$ bpc; the standard metric diluted the phenomenon by $80\times$;
and the efficiency claim inverted without custom kernels while the cache claim
reproduced exactly.

The reproducibility question for frontier architecture research is therefore
not simply one of compute. It is that published mechanisms bundle claims with
very different reproduction requirements, and the bundling is invisible. The
KV-cache argument in CSA is verifiable by anyone with a calculator. The
throughput argument is not verifiable without rebuilding the kernel stack. Both
are presented as efficiency. Distinguishing them, we suggest, would do more for
independent verification of architecture research than any amount of additional
compute at our end.

\paragraph{Code.} The implementation, all run logs, and the measurement scripts
supporting each table are available at
\url{https://github.com/BrettCleary/mini-deepseek-v4}.

\begin{thebibliography}{99}
\small
\bibitem{deepseekv4} DeepSeek-AI. \emph{DeepSeek-V4: Towards Highly Efficient
Million-Token Context Intelligence}. Technical report, 2026.
\bibitem{deepseekv32} DeepSeek-AI. \emph{DeepSeek-V3.2: DeepSeek Sparse
Attention}. Technical report, 2025.
\bibitem{shazeer2019mqa} N.~Shazeer. Fast transformer decoding: one write-head
is all you need. \emph{arXiv:1911.02150}, 2019.
\bibitem{zhang2019rmsnorm} B.~Zhang and R.~Sennrich. Root mean square layer
normalization. \emph{NeurIPS}, 2019.
\bibitem{henry2020qknorm} A.~Henry, P.~Dachapally, S.~Pawar, and Y.~Chen.
Query-key normalization for transformers. \emph{Findings of EMNLP}, 2020.
\bibitem{dehghani2023vit22b} M.~Dehghani et~al. Scaling vision transformers to
22 billion parameters. \emph{ICML}, 2023.
\bibitem{dao2022flashattention} T.~Dao, D.~Fu, S.~Ermon, A.~Rudra, and
C.~R\'{e}. FlashAttention: fast and memory-efficient exact attention with
IO-awareness. \emph{NeurIPS}, 2022.
\bibitem{xiao2024streamingllm} G.~Xiao, Y.~Tian, B.~Chen, S.~Han, and
M.~Lewis. Efficient streaming language models with attention sinks.
\emph{ICLR}, 2024.
\bibitem{hutter2006} M.~Hutter. The human knowledge compression prize, 2006.
\url{http://prize.hutter1.net}
\bibitem{rae2019pg19} J.~W. Rae, A.~Potapenko, S.~M. Jayakumar, and
T.~P. Lillicrap. Compressive transformers for long-range sequence modelling.
\emph{ICLR}, 2020.
\bibitem{loshchilov2019adamw} I.~Loshchilov and F.~Hutter. Decoupled weight
decay regularization. \emph{ICLR}, 2019.
\end{thebibliography}

\end{document}
